\textbf{结论名称：}分组求和法（拆分为等差+等比+常数等）

\textbf{适用条件：}数列通项可拆分为有限个可求和的数列之和，常见拆分如“等差$+$等比”“等差$+$常数”“等比$+$常数”等。

\textbf{变量说明：}设数列 $\{c_n\}$ 满足 $c_n=a_n+b_n+\cdots$，其中 $\{a_n\},\{b_n\}$ 为等差、等比或常数数列，其前$n$项和可分别用已知公式计算。记 $S_n=\sum_{k=1}^{n}c_k$。

\textbf{核心公式：}
\[
\boxed{S_n = \sum_{k=1}^{n} a_k + \sum_{k=1}^{n} b_k + \cdots}
\]
典型拆分：当 $a_n=An+B$（等差），$b_n=C\cdot q^{n-1}$（等比）时，
\[
S_n = \frac{n(A+B+An+B)}{2} + \frac{C(1-q^{n})}{1-q}\quad(q\neq1).
\]

\textbf{使用场景：}通项由一次多项式与指数函数组合构成时，直接分组分别求和，避免整体构造错位相减。

\textbf{简单例子：}
已知 $c_n = 2n + 3^n$，求 $S_n$。
解：拆为 $a_n=2n$（等差），$b_n=3^n$（等比）。
\[
S_n = 2\sum_{k=1}^{n} k + \sum_{k=1}^{n} 3^k = 2\cdot\frac{n(n+1)}{2} + \frac{3(3^{n}-1)}{3-1} = n(n+1) + \frac{3^{n+1}-3}{2}.
\]
验证：$n=1$ 时 $c_1=5$，$S_1=1\cdot2+\frac{9-3}{2}=5$，吻合。